Systems that maintain sorted collections face a recurring design question:
\emph{use a flat sorted array, or a dynamic ordered structure such as a B-tree?}
Arrays offer $O(1)$ indexed access and superior cache locality, but re-sorting
after updates costs $O(n \log n)$. B-trees support $O(\log n)$ updates but
incur $O(n)$ traversal for a full sorted scan or $O(r)$ per rank query.

We study this trade-off under the \emph{update-aware sorting} model, where
the positions of updated elements are explicitly known. Under this model,
\DeltaSort~\cite{deltasort} reduces re-sort cost to $O(n\sqrt{k})$ for $k$ dirty
elements, making sorted arrays competitive with B-trees even at modest update rates.

We make three contributions:
\textbf{(1)} We derive a closed-form crossover formula: for any workload with
$\Qbreak \leq 2$ indexed reads per update epoch, a sorted array with \DeltaSort
outperforms a BTreeMap, regardless of $n$ and $k$.
We validate this analytically against Rust microbenchmarks across
$n \in [10^3, 10^6]$ and $k \in [1, 10^4]$.
\textbf{(2)} We evaluate \DeltaSort on a realistic MIDI sequencer editing workload
($n \leq 10^5$ notes, edit operations modeled on Ardour's sequencer),
showing 5--20$\times$ throughput gains over a \texttt{BTreeSet} for common
drag and quantize operations.
\textbf{(3)} We honestly characterize two domains where \DeltaSort does \emph{not}
provide a practical advantage: large-scale HPC particle sorting ($n = 10^6$, where
\ESM dominates at $k/n \geq 5\%$) and GPU-bound 3D rendering (where
phase-sort cost is $<5\%$ of frame budget regardless of algorithm).
These negative results delineate when the update-aware model is---and is not---the
right abstraction.
